\documentclass[a4paper]{arrowhead}

\usepackage[yyyymmdd]{datetime}
\usepackage{etoolbox}
\usepackage[T1]{fontenc}
\usepackage[utf8]{inputenc}
\usepackage{multirow}
\usepackage{longtable}
\usepackage{amssymb}
\usepackage{float}
\usepackage{tikz}
\usetikzlibrary{positioning,arrows.meta,fit,calc,shapes.misc,shapes.geometric}

\renewcommand{\dateseparator}{-}
\setlength{\emergencystretch}{3em}
\lstset{basicstyle=\ttfamily\footnotesize,breaklines=true,columns=fullflexible,frame=single,rulecolor=\color{gray!50},backgroundcolor=\color{gray!6}}
\tikzset{
  ahsys/.style={draw=ArrowheadBlue,very thick,fill=ArrowheadBlue!6,rounded corners=2pt,minimum width=3.2cm,minimum height=1.1cm,align=center,font=\small\bfseries},
  ahext/.style={draw=ArrowheadGrey,thick,dashed,fill=gray!5,rounded corners=2pt,minimum width=3.2cm,minimum height=1.0cm,align=center,font=\small},
  ahsrv/.style={draw=ArrowheadPurple,thick,fill=ArrowheadPurple!8,rounded corners=6pt,minimum height=0.55cm,inner xsep=4pt,align=center,font=\scriptsize\ttfamily},
  ahcsm/.style={draw=ArrowheadTeal,thick,fill=ArrowheadTeal!8,rounded corners=6pt,minimum height=0.55cm,inner xsep=4pt,align=center,font=\scriptsize\ttfamily},
  ahdb/.style={draw=ArrowheadOrange,thick,fill=ArrowheadOrange!8,cylinder,shape border rotate=90,aspect=0.25,minimum height=0.9cm,align=center,font=\scriptsize},
  ahflow/.style={-{Stealth[length=2.2mm]},thick,ArrowheadGrey},
  ahlabel/.style={font=\scriptsize,fill=white,inner sep=1pt,align=center},
}

\newcommand{\fparam}[1]{\textit{\textcolor{ArrowheadBlue}{#1}}}
%% Special references
\newcommand{\mref}[1]{{\textcolor{ArrowheadPurple}{\hyperref[sec:model:#1]{#1}}}}
\newcommand{\oref}[1]{{\textcolor{ArrowheadBlue}{\hyperref[sec:operations:#1]{#1}}}}
\newcommand{\pdef}[1]{{\textcolor{ArrowheadGrey}{#1\label{sec:model:primitives:#1}\label{sec:model:primitives:#1s}}}}
\newcommand{\pref}[1]{{\textcolor{ArrowheadGrey}{\hyperref[sec:model:primitives:#1]{#1}}}}

\newrobustcmd\osubsection[5]{
  \addtocounter{subsection}{1}
  \addcontentsline{toc}{subsection}{\protect\numberline{\thesubsection}operation \textcolor{ArrowheadBlue}{#1}}
  \subsection*{
    \thesubsection\quad
    #2 \textcolor{ArrowheadPurple}{#3} \\
    \small
    \hspace*{0.075\textwidth}\begin{minipage}{0.1\textwidth}
      \vspace*{1mm}
      Operation: \\
      \notblank{#4}{Input: \\}{}
      \notblank{#5}{Output: \\}{}
    \end{minipage}
    \begin{minipage}{0.825\textwidth}
      \vspace*{1mm}
      \textcolor{ArrowheadBlue}{#1} \\
      \notblank{#4}{\mref{#4} \\}{}
      \notblank{#5}{\mref{#5} \\}{}
    \end{minipage}
  }
  \label{sec:operations:#1}
}
\newrobustcmd\msubsection[2]{
  \addtocounter{subsection}{1}
  \addcontentsline{toc}{subsection}{\protect\numberline{\thesubsection}#1 \textcolor{ArrowheadPurple}{#2}}
  \subsection*{\thesubsection\quad#1 \textcolor{ArrowheadPurple}{#2}}
  \label{sec:model:#2} \label{sec:model:#2s}
}

\begin{document}

%% Arrowhead Document Properties
\ArrowheadTitle{workflowInstance HTTP/TLS/JSON}
\ArrowheadServiceID{workflowInstance}
\ArrowheadType{Interface Design Description}
\ArrowheadTypeShort{IDD}
\ArrowheadVersion{5.4.0}
\ArrowheadDate{2026-11-23}
\ArrowheadAuthor{Documentation maintenance}
\ArrowheadStatus{REVIEW DRAFT}
\ArrowheadContact{Repository maintainers}
\ArrowheadFooter{\href{www.arrowhead.eu}{www.arrowhead.eu}}
\ArrowheadSetup
%%

%% Front Page
\begin{center}
  \vspace*{1cm}
  \huge{\arrowtitle}

  \vspace*{0.2cm}
  \LARGE{\arrowtype}
  \vspace*{1cm}

  \Large{Service ID: \textit{``\arrowid''}}
  \vspace*{\fill}

  

  \vspace*{1cm}
  \vspace*{\fill}

  % Front Page Abstract
  \begin{abstract}
    This document describes the HTTP/TLS/JSON interface of the \texttt{workflowInstance} service of the Eclipse Arrowhead workflow support systems, service version 1.0.0, documentation edition 5.4.0. It realises the abstract operations of the \texttt{workflowInstance} SD.
  \end{abstract}

  \vspace*{1cm}
\end{center}

\newpage
%%

%% Table of Contents
\tableofcontents
\newpage
%%

\section{Overview}
\label{sec:overview}
This document describes how the \texttt{workflowInstance} service, which starts, observes and aborts workflow sessions, is realised with HTTP~\cite{rfc9110}
over TLS~\cite{rfc8446} and JSON~\cite{rfc8259} payloads. Request and response bodies use content type
\texttt{application/json}; JSON field names are case-sensitive and fields not listed in this document must not be sent
(the recipe decoder rejects unknown fields). Optional fields may be omitted; absent and \texttt{null} are equivalent.

The service is offered by the WorkflowManager on its Arrowhead connector (default port 8482) with base path \texttt{/\allowbreak{}workflowmanager/\allowbreak{}instance}.
The ServiceRegistry entry lists every operation with its HTTP method and path relative to the base path.

The rest of this document is organized as follows. Section~\ref{sec:operations} specifies each operation,
Section~\ref{sec:model} the JSON data models, Section~\ref{sec:security} the security requirements,
Section~\ref{sec:errors} the error representation and Section~\ref{sec:implementations} the implementations.

\newpage
\section{Service Operations}
\label{sec:operations}

{\small
\begin{longtable}{|p{3.0cm}|p{1.5cm}|p{5.3cm}|p{3.0cm}|p{3.0cm}|}

\hline
\rowcolor{gray!33} Operation & Method & Path & Input & Output \\ \hline
\endhead
\oref{start} & POST & \texttt{/\allowbreak{}workflowmanager/\allowbreak{}instance/\allowbreak{}start} & \mref{StartSessionRequest} & \mref{Session} \\ \hline
\oref{abort} & DELETE & \texttt{/\allowbreak{}workflowmanager/\allowbreak{}instance/\allowbreak{}abort/\allowbreak{}\{id\}} & \mref{SessionRef} & \mref{Session} \\ \hline
\oref{query} & GET & \texttt{/\allowbreak{}workflowmanager/\allowbreak{}instance/\allowbreak{}sessions} & \mref{SessionFilter} & \mref{SessionList} \\ \hline
\oref{get-status} & GET & \texttt{/\allowbreak{}workflowmanager/\allowbreak{}instance/\allowbreak{}sessions/\allowbreak{}\{id\}} & \mref{SessionRef} & \mref{Session} \\ \hline
\oref{get-trace} & GET & \texttt{/\allowbreak{}workflowmanager/\allowbreak{}instance/\allowbreak{}sessions/\allowbreak{}\{id\}/\allowbreak{}trace} & \mref{SessionRef} & \mref{Trace} \\ \hline
\end{longtable}
}

\osubsection{start}{POST}{/workflowmanager/instance/start}{StartSessionRequest}{Session}
Creates a session of an active recipe version (or of the product's recipe), fires the first transitions and dispatches the enabled steps.

{\small
\begin{longtable}{|p{2.0cm}|p{13.8cm}|}

\hline
\rowcolor{gray!33} Status & When \\ \hline
\endhead
Success & 201 Created with the new session. \\ \hline
400 & Missing request or recipe, parameters not an object, required parameter missing, unknown dispatch mode, or a product recipe that does not validate. \\ \hline
404 & The referenced object does not exist (\texttt{DATA\_NOT\_FOUND}). \\ \hline
409 & The recipe version is retired. \\ \hline
503 & The product's \texttt{workflowProduct} returned no recipe (\texttt{EXTERNAL\_SERVER\_ERROR}). \\ \hline
401 & No valid identity token, or the token is bound to another certificate (\texttt{AUTH}). \\ \hline
403 & The consumer has no ConsumerAuthorization grant for this operation (\texttt{FORBIDDEN}). \\ \hline
\end{longtable}
}


\begin{lstlisting}[caption={Example request body}]
{ "recipe": "temperature-branching", "version": "1.0.0",
  "subject": "period-2026-11-21T10:00", "parameters": { "threshold": 23 },
  "dispatch": "push" }
\end{lstlisting}


\begin{lstlisting}[caption={Example response body}]
{ "id": 42, "recipe": "temperature-branching", "version": "1.0.0",
  "subject": "period-2026-11-21T10:00", "owner": "WfClient", "status": "running",
  "dispatch": "push", "extensions": [],
  "stepMarking": { "read": "running", "evaluate": "inactive", "heatOn": "inactive",
                   "heatOff": "inactive", "confirm": "inactive" },
  "actionMarking": { "control": "active" }, "modelTransitions": 2,
  "transitionBound": 17, "createdAt": "...", "updatedAt": "..." }
\end{lstlisting}

\osubsection{abort}{DELETE}{/workflowmanager/instance/abort/\{id\}}{SessionRef}{Session}
Aborts a running session (extension x-abort). Its active work items become \texttt{ABORTED} and their executors receive \texttt{workflowExecution.abort}; leases are released when \texttt{workflow.lease.release-on-failure} is set.

{\small
\begin{longtable}{|p{3.0cm}|p{1.4cm}|p{3.2cm}|p{8.2cm}|}

\hline
\rowcolor{gray!33} Parameter & In & Type & Description \\ \hline
\endhead
\fparam{id} & path & \pref{SessionId} & Session. \\ \hline
\fparam{reason} & query & \pref{String} & Optional reason. \\ \hline
\end{longtable}
}


{\small
\begin{longtable}{|p{2.0cm}|p{13.8cm}|}

\hline
\rowcolor{gray!33} Status & When \\ \hline
\endhead
Success & 200 OK with the failed session. \\ \hline
403 & The caller is neither the owner nor a sysop. \\ \hline
404 & The referenced object does not exist (\texttt{DATA\_NOT\_FOUND}). \\ \hline
409 & The session is already completed or failed. \\ \hline
401 & No valid identity token, or the token is bound to another certificate (\texttt{AUTH}). \\ \hline
403 & The consumer has no ConsumerAuthorization grant for this operation (\texttt{FORBIDDEN}). \\ \hline
\end{longtable}
}

\osubsection{query}{GET}{/workflowmanager/instance/sessions}{SessionFilter}{SessionList}
Lists sessions, newest first, optionally in one status.

{\small
\begin{longtable}{|p{3.0cm}|p{1.4cm}|p{3.2cm}|p{8.2cm}|}

\hline
\rowcolor{gray!33} Parameter & In & Type & Description \\ \hline
\endhead
\fparam{status} & query & \mref{SessionStatus} & Optional status. \\ \hline
\end{longtable}
}


{\small
\begin{longtable}{|p{2.0cm}|p{13.8cm}|}

\hline
\rowcolor{gray!33} Status & When \\ \hline
\endhead
Success & 200 OK. \\ \hline
400 & Unknown status. \\ \hline
401 & No valid identity token, or the token is bound to another certificate (\texttt{AUTH}). \\ \hline
403 & The consumer has no ConsumerAuthorization grant for this operation (\texttt{FORBIDDEN}). \\ \hline
\end{longtable}
}

\osubsection{get-status}{GET}{/workflowmanager/instance/sessions/\{id\}}{SessionRef}{Session}
Returns the markings and status of one session.

{\small
\begin{longtable}{|p{3.0cm}|p{1.4cm}|p{3.2cm}|p{8.2cm}|}

\hline
\rowcolor{gray!33} Parameter & In & Type & Description \\ \hline
\endhead
\fparam{id} & path & \pref{SessionId} & Session. \\ \hline
\end{longtable}
}


{\small
\begin{longtable}{|p{2.0cm}|p{13.8cm}|}

\hline
\rowcolor{gray!33} Status & When \\ \hline
\endhead
Success & 200 OK. \\ \hline
404 & The referenced object does not exist (\texttt{DATA\_NOT\_FOUND}). \\ \hline
401 & No valid identity token, or the token is bound to another certificate (\texttt{AUTH}). \\ \hline
403 & The consumer has no ConsumerAuthorization grant for this operation (\texttt{FORBIDDEN}). \\ \hline
\end{longtable}
}

\osubsection{get-trace}{GET}{/workflowmanager/instance/sessions/\{id\}/trace}{SessionRef}{Trace}
Returns the run: every transition, $\xi$, $\beta$, $\varepsilon$, and whether the run stayed within the proven model.

{\small
\begin{longtable}{|p{3.0cm}|p{1.4cm}|p{3.2cm}|p{8.2cm}|}

\hline
\rowcolor{gray!33} Parameter & In & Type & Description \\ \hline
\endhead
\fparam{id} & path & \pref{SessionId} & Session. \\ \hline
\end{longtable}
}


{\small
\begin{longtable}{|p{2.0cm}|p{13.8cm}|}

\hline
\rowcolor{gray!33} Status & When \\ \hline
\endhead
Success & 200 OK. \\ \hline
404 & The referenced object does not exist (\texttt{DATA\_NOT\_FOUND}). \\ \hline
401 & No valid identity token, or the token is bound to another certificate (\texttt{AUTH}). \\ \hline
403 & The consumer has no ConsumerAuthorization grant for this operation (\texttt{FORBIDDEN}). \\ \hline
\end{longtable}
}


\begin{lstlisting}[caption={Example response body}]
{ "sessionId": 42, "status": "completed",
  "trace": [ { "seq": 1, "rule": "S-Enable", "target": "read", "at": "..." },
             { "seq": 2, "rule": "S-Bind", "target": "read", "detail": "WfeA", "at": "..." },
             ... ,
             { "seq": 14, "rule": "Z-Close", "at": "..." } ],
  "results": { "read": { "temperature": 21.5 }, "evaluate": { "heat": true }, ... },
  "bindings": { "read": { "provider": "WfEquipment", ... } },
  "executorAssignment": { "read": "WfeA", ... }, "runner": { "read": "WfeA", ... },
  "modelTransitions": 14, "transitionBound": 17, "withinModel": false }
\end{lstlisting}


\newpage
\section{Data Models}
\label{sec:model}
The JSON objects exchanged by the operations, in alphabetic order. Field types refer to other data models or to the
primitives of Section~\ref{sec:model:primitives}.

\msubsection{enum}{DispatchMode}
How a step reaches an executor.
 Encoded as a JSON string.
{\small
\begin{longtable}{|p{3.9cm}|p{11.9cm}|}

\hline
\rowcolor{gray!33} Value & Meaning \\ \hline
\endhead
\texttt{push} & The WorkflowManager selects an executor and calls \texttt{workflowExecution.start} (default). \\ \hline
\texttt{pull} & The WorkflowManager publishes the work item; an executor claims it with \texttt{workflowWorkItem.claim}. \\ \hline
\texttt{event} & Reserved for EventHandler notification; currently handled as \texttt{pull}. \\ \hline
\end{longtable}
}

\msubsection{struct}{ErrorMessage}
Error response of every operation (common library \texttt{ErrorMessageDTO}).

{\small
\begin{longtable}{|p{3.7cm}|p{4.2cm}|p{1.6cm}|p{6.3cm}|}

\hline
\rowcolor{gray!33} Field & Type & Mandatory & Description \\ \hline
\endhead
\texttt{errorMessage} & \pref{String} & yes & Explanation. \\ \hline
\texttt{errorCode} & \pref{Integer} & yes & HTTP status code. \\ \hline
\texttt{exceptionType} & \pref{String} & yes & \texttt{INVALID\_PARAMETER}, \texttt{AUTH}, \texttt{FORBIDDEN}, \texttt{DATA\_NOT\_FOUND}, \texttt{EXTERNAL\_SERVER\_ERROR}, \ldots \\ \hline
\texttt{origin} & \pref{String} & no & The operation that failed, e.g. \texttt{POST /workflowmanager/work-item/report}. \\ \hline
\end{longtable}
}

\msubsection{enum}{Rule}
Transition rules: the eight rules of~\cite{varga2026} and the extensions (\texttt{x-}).
 Encoded as a JSON string.
{\small
\begin{longtable}{|p{3.9cm}|p{11.9cm}|}

\hline
\rowcolor{gray!33} Value & Meaning \\ \hline
\endhead
\texttt{A-Start} & An action becomes active after its predecessor completed. \\ \hline
\texttt{S-Enable} & A step becomes enabled when its action is active and its predecessors completed. \\ \hline
\texttt{S-Bind} & An enabled step is bound to a provider (and executor) and becomes running. \\ \hline
\texttt{S-Done} & A running step completed; its result is stored. \\ \hline
\texttt{A-Finish} & Every step of an action completed. \\ \hline
\texttt{Z-Close} & Every action completed: the session is completed. \\ \hline
\texttt{S-FailBind} & No executor/provider could be bound: the session fails. \\ \hline
\texttt{S-FailInvoke} & The invocation failed: the session fails. \\ \hline
\texttt{x-rebind} & Extension: the step is bound again after \texttt{ERROR} or an unreachable executor. \\ \hline
\texttt{x-abort} & Extension: the session was aborted. \\ \hline
\texttt{x-timeout} & Extension: a step exceeded its timeout. \\ \hline
\texttt{x-skip} & Extension: a step was skipped because no guarded transition to it was taken. \\ \hline
\end{longtable}
}

\msubsection{struct}{Session}
The state $z = (M, M_A, \beta, \varepsilon, \xi, status)$ of a session, as seen from outside.

{\small
\begin{longtable}{|p{3.7cm}|p{4.2cm}|p{1.6cm}|p{6.3cm}|}

\hline
\rowcolor{gray!33} Field & Type & Mandatory & Description \\ \hline
\endhead
\texttt{id} & \pref{SessionId} & yes & Session identifier. \\ \hline
\texttt{recipe} & \pref{Name} & yes & Recipe name. \\ \hline
\texttt{version} & \pref{Version} & yes & Recipe version. \\ \hline
\texttt{subject} & \pref{String} & no & Subject of the session. \\ \hline
\texttt{owner} & \pref{SystemName} & yes & The system that started the session; it may abort it. \\ \hline
\texttt{status} & \mref{SessionStatus} & yes & \texttt{running}, \texttt{completed} or \texttt{failed}. \\ \hline
\texttt{dispatch} & \mref{DispatchMode} & yes & Effective dispatch mode. \\ \hline
\texttt{failure} & \pref{String} & no & Why the session failed. \\ \hline
\texttt{extensions} & \pref{List}$<$\pref{String}$>$ & yes & Extensions used (\texttt{x-rebind}, \texttt{x-abort}, \texttt{x-timeout}, \texttt{x-skip}, \ldots); empty when the run stays inside the proven model. \\ \hline
\texttt{stepMarking} & \pref{Map}$<$\pref{Name}, \pref{String}$>$ & yes & $M$: \texttt{inactive}, \texttt{enabled}, \texttt{running}, \texttt{completed} or \texttt{skipped} per step. \\ \hline
\texttt{actionMarking} & \pref{Map}$<$\pref{Name}, \pref{String}$>$ & yes & $M_A$: \texttt{inactive}, \texttt{active} or \texttt{completed} per action. \\ \hline
\texttt{modelTransitions} & \pref{Integer} & yes & Number of model transitions so far. \\ \hline
\texttt{transitionBound} & \pref{Integer} & yes & Theorem~1 bound $3|S| + 2|A|$. \\ \hline
\texttt{product} & \pref{SystemName} & no & Product system, when the session came from one. \\ \hline
\texttt{createdAt} & \pref{DateTime} & yes & Start time. \\ \hline
\texttt{updatedAt} & \pref{DateTime} & yes & Time of the last transition. \\ \hline
\end{longtable}
}

\msubsection{struct}{SessionFilter}
Optional filter of \oref{query}.

{\small
\begin{longtable}{|p{3.7cm}|p{4.2cm}|p{1.6cm}|p{6.3cm}|}

\hline
\rowcolor{gray!33} Field & Type & Mandatory & Description \\ \hline
\endhead
\texttt{status} & \mref{SessionStatus} & no & Only sessions in this status. \\ \hline
\end{longtable}
}

\msubsection{struct}{SessionList}
Sessions, newest first.

{\small
\begin{longtable}{|p{3.7cm}|p{4.2cm}|p{1.6cm}|p{6.3cm}|}

\hline
\rowcolor{gray!33} Field & Type & Mandatory & Description \\ \hline
\endhead
\texttt{entries} & \pref{List}$<$\mref{Session}$>$ & yes & The sessions. \\ \hline
\texttt{count} & \pref{Integer} & yes & Number of entries. \\ \hline
\end{longtable}
}

\msubsection{struct}{SessionRef}
Identifies a session; \oref{abort} also takes an optional reason.

{\small
\begin{longtable}{|p{3.7cm}|p{4.2cm}|p{1.6cm}|p{6.3cm}|}

\hline
\rowcolor{gray!33} Field & Type & Mandatory & Description \\ \hline
\endhead
\texttt{id} & \pref{SessionId} & yes & Session identifier. \\ \hline
\texttt{reason} & \pref{String} & no & Reason recorded with an abort. \\ \hline
\end{longtable}
}

\msubsection{enum}{SessionStatus}
Status of a session. \texttt{completed} and \texttt{failed} are final.
 Encoded as a JSON string.
{\small
\begin{longtable}{|p{3.9cm}|p{11.9cm}|}

\hline
\rowcolor{gray!33} Value & Meaning \\ \hline
\endhead
\texttt{running} & Some rule may still apply. \\ \hline
\texttt{completed} & Z-Close fired: every action completed. \\ \hline
\texttt{failed} & S-FailBind, S-FailInvoke or an extension (abort, timeout) ended the session. \\ \hline
\end{longtable}
}

\msubsection{struct}{StartSessionRequest}
Starts one session (one execution of one recipe version for one subject).

{\small
\begin{longtable}{|p{3.7cm}|p{4.2cm}|p{1.6cm}|p{6.3cm}|}

\hline
\rowcolor{gray!33} Field & Type & Mandatory & Description \\ \hline
\endhead
\texttt{recipe} & \pref{Name} & no & Recipe name; may be omitted when \texttt{product} is given. \\ \hline
\texttt{version} & \pref{Version} & no & Recipe version; the latest active version when absent. \\ \hline
\texttt{subject} & \pref{String} & no & What the session works on: product unit, control period or batch. \\ \hline
\texttt{parameters} & \pref{JsonObject} & no & Session parameters $\rho$; required parameters must be present. \\ \hline
\texttt{dispatch} & \mref{DispatchMode} & no & Overrides the recipe's dispatch mode. \\ \hline
\texttt{product} & \pref{SystemName} & no & A product system providing \texttt{workflowProduct}: the recipe and parameters are retrieved from it, and the results are stored back at Z-Close. \\ \hline
\end{longtable}
}

\msubsection{struct}{Trace}
The run of a session: the transitions and the data the rules produced.

{\small
\begin{longtable}{|p{3.7cm}|p{4.2cm}|p{1.6cm}|p{6.3cm}|}

\hline
\rowcolor{gray!33} Field & Type & Mandatory & Description \\ \hline
\endhead
\texttt{sessionId} & \pref{SessionId} & yes & Session identifier. \\ \hline
\texttt{status} & \mref{SessionStatus} & yes & Current status. \\ \hline
\texttt{trace} & \pref{List}$<$\mref{Transition}$>$ & yes & Transitions in order. \\ \hline
\texttt{results} & \pref{Map}$<$\pref{Name}, \pref{JsonValue}$>$ & yes & $\xi$: stored step results. \\ \hline
\texttt{bindings} & \pref{Map}$<$\pref{Name}, \pref{JsonValue}$>$ & yes & $\beta$: orchestration result per bound step. \\ \hline
\texttt{executorAssignment} & \pref{Map}$<$\pref{Name}, \pref{SystemName}$>$ & yes & $\varepsilon$: executor per step. \\ \hline
\texttt{runner} & \pref{Map}$<$\pref{Name}, \pref{SystemName}$>$ & yes & The system that ran each step. \\ \hline
\texttt{modelTransitions} & \pref{Integer} & yes & Transitions of the proven rule set. \\ \hline
\texttt{transitionBound} & \pref{Integer} & yes & Theorem~1 bound. \\ \hline
\texttt{withinModel} & \pref{Boolean} & yes & True when no extension was used, so Theorems~1--2 apply. \\ \hline
\end{longtable}
}

\msubsection{struct}{Transition}
One transition of the session.

{\small
\begin{longtable}{|p{3.7cm}|p{4.2cm}|p{1.6cm}|p{6.3cm}|}

\hline
\rowcolor{gray!33} Field & Type & Mandatory & Description \\ \hline
\endhead
\texttt{seq} & \pref{Integer} & yes & Sequence number from 1. \\ \hline
\texttt{rule} & \mref{Rule} & yes & The rule that fired. \\ \hline
\texttt{target} & \pref{Name} & no & The step or action it applied to. \\ \hline
\texttt{detail} & \pref{String} & no & Additional information, e.g. the executor or failure. \\ \hline
\texttt{at} & \pref{DateTime} & yes & When it fired. \\ \hline
\end{longtable}
}


\subsection{Primitives}
\label{sec:model:primitives}
The abstract primitive types of the SD are represented in JSON~\cite{rfc8259} as follows.
\subsubsection*{alias \pdef{Boolean}}
JSON \texttt{true} or \texttt{false}.

\subsubsection*{alias \pdef{DateTime}}
JSON string in the RFC 3339 format~\cite{rfc3339}, e.g. \texttt{2026-11-23T10:15:30Z}.

\subsubsection*{alias \pdef{Integer}}
JSON number without fraction.

\subsubsection*{alias \pdef{JsonObject}}
JSON object.

\subsubsection*{alias \pdef{JsonValue}}
Any JSON value~\cite{rfc8259}.

\subsubsection*{alias \pdef{List}}
JSON array.

\subsubsection*{alias \pdef{Map}}
JSON object whose member names are the keys.

\subsubsection*{alias \pdef{Name}}
JSON string matching \texttt{\^{}[A-Za-z][A-Za-z0-9\_-]\{0,62\}\$}.

\subsubsection*{alias \pdef{SessionId}}
JSON number (64-bit integer); a path segment when used in a path.

\subsubsection*{alias \pdef{String}}
JSON string (UTF-8).

\subsubsection*{alias \pdef{SystemName}}
JSON string.

\subsubsection*{alias \pdef{Version}}
JSON string matching \texttt{\^{}\textbackslash{}d+\textbackslash{}.\textbackslash{}d+\textbackslash{}.\textbackslash{}d+\$}.



\section{Security}
\label{sec:security}
The interface is offered only over TLS~\cite{rfc8446} (TLS 1.3 or 1.2) with mutual authentication: the consumer must
present an Arrowhead certificate issued under the local cloud's certificate authority, and certificates on the
cloud's CRL are refused. The service is registered with interface template \texttt{generic\_https} and interface
policy \texttt{CERT\_AUTH}.

Every request must carry the consumer's identity token, obtained from the Authentication system and bound to the
consumer's certificate:
\begin{lstlisting}
Authorization: Bearer IDENTITY-TOKEN//<identity token>
\end{lstlisting}
A token issued to another system's certificate is refused (401). The provider then checks a ConsumerAuthorization
grant for the operation (\texttt{authorization.verify}, \texttt{service.authorization.enabled=true}); a missing grant is
refused with 403. Each operation is authorized separately. \texttt{abort} additionally requires the caller to be the session owner or a sysop.


\section{Errors}
\label{sec:errors}
Every error response carries an \mref{ErrorMessage} body with content type \texttt{application/json}, for example:
\begin{lstlisting}
{ "errorMessage": "Work item 5d2c0b1e-... is held by another executor",
  "errorCode": 403, "exceptionType": "FORBIDDEN",
  "origin": "POST /workflowmanager/work-item/report" }
\end{lstlisting}
Status codes follow HTTP semantics~\cite{rfc9110}: 400 invalid parameter, 401 authentication, 403 forbidden, 404 not
found, 409 conflict with the current state, 500 internal error and 503 failure of an external system.


\section{Implementations}
\label{sec:implementations}
\begin{itemize}
\item \textbf{Java (reference).} \texttt{Eclipse-Arrowhead-java/ah-Workflow-java}, module \texttt{workflow-manager}. Paths, methods and service
operations are defined once in \texttt{workflow-model} (\texttt{wire/WorkflowServices.java}, \texttt{wire/ServiceOperations.java}) and checked by the provider-side authorization filter.
\item \textbf{Go.} Not implemented: \texttt{Eclipse-Arrowhead-go} has no workflow systems. A Go implementation must
provide exactly the paths, methods, JSON field names and status codes of this IDD to be interoperable.
\item \textbf{Rust.} Not implemented: \texttt{Eclipse-Arrowhead-rust} has no workflow systems. The same requirement applies.
\end{itemize}
Until the Go and Rust implementations exist, the workflow systems are \emph{not} interoperable across the three
implementations; only the Java systems interoperate with Java, Go and Rust core systems through the common Arrowhead
interfaces (ServiceRegistry, Authentication, ConsumerAuthorization, DynamicServiceOrchestration).


\newpage
\bibliographystyle{IEEEtran}
\bibliography{bibliography}

\newpage

\section{Revision History}
\subsection{Amendments}

\noindent\begin{tabularx}{\textwidth}{| p{1cm} | p{2.2cm} | p{1.6cm} | X | p{3.6cm} |} \hline
\rowcolor{gray!33} No. & Date & Version & Subject of Amendments & Author \\ \hline
1 & 2026-11-23 & \arrowversion & Initial edition: interface design description of workflowInstance over HTTP/TLS/JSON, written from the Java reference implementation (ah-Workflow-java). & Documentation maintenance \\ \hline
\end{tabularx}

\subsection{Quality Assurance}

\noindent\begin{tabularx}{\textwidth}{| p{1cm} | p{2.2cm} | p{1.6cm} | X |} \hline
\rowcolor{gray!33} No. & Date & Version & Approved by \\ \hline
1 & & \arrowversion & Pending review \\ \hline
\end{tabularx}

\end{document}
